\section*{Supplementary Material A: Edge-Weighting Technical Specification}

\subsection{Edge Weight Formula}

For an edge $(i,j)$ connecting adjacent census tracts $i$ and $j$, the weight is computed as:

\begin{equation}
w(i,j) = \frac{\alpha \cdot s_{demo}(i,j)}{d(i,j) + \epsilon}
\end{equation}

where:
\begin{itemize}
\item $d(i,j)$ = Euclidean distance between tract centroids (kilometers)
\item $s_{demo}(i,j)$ = Demographic similarity score $\in [0,1]$
\item $\alpha$ = Scaling parameter (default: 10.0)
\item $\epsilon$ = Small constant to prevent division by zero (default: 0.1 km)
\end{itemize}

\subsection{Demographic Similarity Measure}

Demographic similarity uses cosine similarity of racial/ethnic composition vectors:

\begin{equation}
s_{demo}(i,j) = \frac{\vec{v}_i \cdot \vec{v}_j}{|\vec{v}_i| \cdot |\vec{v}_j|}
\end{equation}

where $\vec{v}_i = [p_{white}, p_{black}, p_{hispanic}, p_{asian}, p_{other}]$ is the 5-dimensional vector of population proportions for tract $i$, normalized to sum to 1.

\textbf{Rationale}: Cosine similarity ranges from 0 (orthogonal, maximally dissimilar) to 1 (identical proportions). This metric:
\begin{itemize}
\item Treats all racial/ethnic groups symmetrically
\item Is insensitive to absolute population sizes (uses proportions)
\item Produces intuitive results: two tracts both 70\% Black have $s_{demo} \approx 1$
\item Avoids requiring explicit thresholds for ``minority'' vs. ``majority''
\end{itemize}

\subsection{Distance Normalization}

The denominator $d(i,j) + \epsilon$ ensures:
\begin{itemize}
\item Weights decay with distance (inversely proportional)
\item Adjacent tracts (small $d$) receive higher weights than distant tracts
\item The constant $\epsilon = 0.1$ km prevents infinite weights for tracts sharing long boundaries (where centroids may be very close)
\item Typical tract centroid distances: 0.5-5 km within urban areas, 5-50 km in rural areas
\end{itemize}

\subsection{Scaling Parameter $\alpha$}

The parameter $\alpha$ controls the relative importance of demographic similarity versus geographic compactness:
\begin{itemize}
\item $\alpha = 0$: Pure geographic compactness (all weights = 0, METIS minimizes geometric edge-cut)
\item $\alpha = 10$ (default): Balanced weighting for VRA compliance without over-optimization
\item $\alpha = 100$: Strong preference for preserving minority communities (may reduce compactness)
\end{itemize}

In our analysis, $\alpha = 10$ was selected through sensitivity testing to achieve:
\begin{itemize}
\item Sufficient minority-minority district creation to satisfy VRA geographic compactness requirements
\item Maintained overall compactness (Polsby-Popper scores $> 0.3$)
\item Avoided pathological cases where high weights create non-compact districts
\end{itemize}

\subsection{Weight Distribution Statistics}

Across all 50 states and 85,000 census tracts (2020 data), edge weights exhibit the following distribution:

\begin{table}[h]
\centering
\caption{Edge Weight Distribution Statistics (National, 2020)}
\begin{tabular}{lr}
\toprule
\textbf{Statistic} & \textbf{Value} \\
\midrule
Minimum weight & 0.02 \\
5th percentile & 0.15 \\
25th percentile (Q1) & 0.48 \\
Median (Q2) & 1.23 \\
75th percentile (Q3) & 3.76 \\
95th percentile & 12.45 \\
Maximum weight & 48.32 \\
Mean & 2.14 \\
Standard deviation & 3.89 \\
Coefficient of variation & 1.82 \\
\midrule
\multicolumn{2}{l}{\textit{Skewness and heavy tails}} \\
Skewness & 4.23 \\
\% weights $> 10$ & 6.2\% \\
\% weights $> 20$ & 1.4\% \\
\bottomrule
\end{tabular}
\end{table}

\textbf{Interpretation}:
\begin{itemize}
\item \textbf{Well-behaved distribution}: Median weight (1.23) is moderate, with most edges having weights 0.5-4.0
\item \textbf{No pathological cases}: Maximum weight (48.32) occurs for highly similar, adjacent minority tracts in dense urban areas. This is geometrically sensible: such tracts should strongly prefer staying together
\item \textbf{Positive skewness}: Right tail (high weights) reflects intentional design: demographically similar tracts receive substantially higher weights, enabling VRA compliance
\item \textbf{Low weights dominate}: 75\% of edges have weights $< 3.76$, ensuring geographic compactness remains primary optimization objective
\end{itemize}

\subsection{Regional Variation}

Weight distributions vary by urban/rural geography:

\begin{table}[h]
\centering
\caption{Edge Weight Statistics by Region Type}
\begin{tabular}{lrrr}
\toprule
\textbf{Statistic} & \textbf{Urban} & \textbf{Suburban} & \textbf{Rural} \\
\midrule
Median weight & 1.45 & 1.18 & 0.89 \\
95th percentile & 15.23 & 11.87 & 8.34 \\
\% weights $> 10$ & 8.1\% & 5.9\% & 4.2\% \\
\bottomrule
\end{tabular}
\end{table}

Urban areas show higher weights due to:
\begin{itemize}
\item Greater demographic diversity (more variance in $s_{demo}$)
\item Shorter inter-tract distances (smaller $d(i,j)$)
\item Concentrated minority populations (higher $s_{demo}$ for minority-minority edges)
\end{itemize}

\subsection{Sensitivity to $\alpha$}

To verify weights are not overly sensitive to $\alpha$ choice, we tested alternative scaling parameters:

\begin{table}[h]
\centering
\caption{Sensitivity to Scaling Parameter $\alpha$}
\begin{tabular}{lrrr}
\toprule
\textbf{$\alpha$} & \textbf{Median Weight} & \textbf{MM Districts} & \textbf{Compactness (PP)} \\
\midrule
1.0 & 0.12 & 98 & 0.42 \\
5.0 & 0.62 & 124 & 0.38 \\
10.0 (default) & 1.23 & 137 & 0.35 \\
20.0 & 2.46 & 143 & 0.31 \\
50.0 & 6.15 & 149 & 0.26 \\
\bottomrule
\end{tabular}
\label{tab:alpha-sensitivity}
\end{table}

Results show:
\begin{itemize}
\item $\alpha = 10$ balances VRA compliance (137 MM districts) with compactness (0.35 Polsby-Popper)
\item Doubling $\alpha$ to 20 yields diminishing returns (+6 MM districts) with compactness penalty
\item $\alpha = 50$ produces excessive VRA focus at cost of non-compact districts
\end{itemize}

\subsection{Computational Considerations}

Edge weights are computed once during graph construction, before METIS partitioning. For 85,000 tracts:
\begin{itemize}
\item Adjacency detection: 2.3 million edges (queen contiguity)
\item Weight computation: $< 1$ second (vectorized operations)
\item Memory footprint: 37 MB (sparse edge list representation)
\item METIS optimization time: 5-120 seconds depending on state size
\end{itemize}

Weights remain static during METIS refinement passes---they are not updated as partitions evolve. This ensures reproducibility: identical inputs always produce identical outputs.

\subsection{Comparison to Alternative Weighting Schemes}

We tested three alternative demographic similarity measures:

\begin{enumerate}
\item \textbf{Binary minority indicator}: $s_{demo} = 1$ if both tracts $> 50\%$ minority, else 0
\item \textbf{Hamming distance}: Count of demographic categories with $> 10\%$ difference
\item \textbf{Euclidean distance}: $1 - |\vec{v}_i - \vec{v}_j|_2 / \sqrt{2}$
\end{enumerate}

Cosine similarity outperforms alternatives because:
\begin{itemize}
\item Binary indicator is too coarse (only captures 50\% threshold, ignores gradations)
\item Hamming distance requires arbitrary thresholds (10\%? 20\%?)
\item Euclidean distance over-penalizes small differences in large categories
\item Cosine similarity is threshold-free and proportionally weighted
\end{itemize}

\subsection{Summary}

The edge-weighting scheme $w(i,j) = \alpha \cdot s_{demo}(i,j) / (d(i,j) + \epsilon)$ provides:
\begin{itemize}
\item \textbf{Mathematically precise}: Clear formula with defined parameters
\item \textbf{Well-behaved}: Weight distribution shows no pathological cases
\item \textbf{Tunable}: Parameter $\alpha$ allows balancing VRA vs. compactness
\item \textbf{Transparent}: All inputs (demographics, distance) are public census data
\item \textbf{Reproducible}: Static weights + deterministic METIS = reproducible outputs
\end{itemize}

This specification enables independent verification and replication of our results.
